\begin{figure}
\centering
\captionsetup[subfigure]{font=footnotesize,skip=0pt}

\begin{subfigure}[t]{0.696\linewidth}
\centering
\includegraphics[width=\linewidth, trim=0.8cm 0 0 0, clip]{../binder/binder-2025-11-30-external-selfbb-truebb.ipynb/binder/teeplots/2025-11-30-external-selfbb-truebb/hue=kind1+style=kind1+viz=lineplot+x=competition-stint+y=focal-prevalence+ext=.pdf}
\caption{transfer rounds 0 to 100}
\label{fig:fitness-background:timeseries}
\end{subfigure}%
\begin{subfigure}[t]{0.304\linewidth}
\centering
\includegraphics[width=\linewidth, trim=1.3cm 0 0 0.1cm, clip]{../binder/binder-2025-11-30-external-selfbb-truebb.ipynb/binder/teeplots/2025-11-30-external-selfbb-truebb/hue=kind+inner=quartile+viz=violinplot+x=kind+y=focal-prevalence+ext=.pdf}
\caption{rounds 15 to 100}
\label{fig:fitness-background:distribution}
\end{subfigure}

\caption{%
\textbf{Patterned growth creates biogenic niche.}
\footnotesize
Around transfer round 15, fitness of case study strain in self-context diverges significantly above other contexts.
For background type ``None,'' competition is initialized 50\%/50\% between compared strains (Figure \ref{fig:fit:nobbg}).
Otherwise, competition is initialized 25\%/25\%/50\% between compared strains and background (Figure \ref{fig:fit:bbg}).
Shaded bands are bootstrap 95\% CI.
Mann-Whitney test with Cliff's $\delta$ reported.
% From this point onwards, self context affords significantly higher tournament fitness than ecological context (Mann-Whitney $U$ test, $p<10^{-25}$).
% By contrast, ecological context more closely tracks background-free tournament fitness, though exhibiting strong episodic deviations.
% Absent background context, the case study strain's tournament fitness fell below null expectation of 0.5 ($\mu=0.39$; $[0.31,\;0.43]$ bootstrap 95\% CI, $n=40$).
% By contrast, in self-context tournament fitness fell above null expectation ($\mu=0.72$; $[0.69,\;0.76]$ bootstrap 95\% CI, $n=40$).
}
\label{fig:fitness-background}
\end{figure}
